\documentclass[10pt,a4paper]{amsart}
\usepackage[margin=19mm]{geometry}
\usepackage{amsmath,amssymb,mathtools}
\usepackage[scr=rsfs,cal=euler]{mathalfa}
\usepackage{xfrac}
\usepackage[unicode,hidelinks,bookmarks=false]{hyperref}
\newcommand{\ie}{i.e.\ }
\newcommand{\cf}{cf.\ }
\newcommand{\resp}{respectively\ }
\newcommand{\Subsection}{\subsection}
\providecommand{\nobreakdash}{\nobreak\hbox{-}}
\setlength{\emergencystretch}{2em}
\multlinegap0pt
\usepackage[shortlabels]{enumitem}
\usepackage{amssymb}
\usepackage{tikz}
\usepackage{algorithm}\usepackage{algpseudocode}
\usepackage{float}
\newtheorem{lemma}{Lemma}
\newcommand{\EE}{\mathbb{E}}
\title{Planar 1-ended graphs can be periodically coloured}
\author{Luke Waite}
\date{Source: arXiv:2604.22705v1 (CC BY 4.0)}
\begin{document}
\maketitle

\noindent\emph{Source and licence.} Planar 1-ended graphs can be periodically coloured. Version arXiv:2604.22705v1 (CC BY 4.0), distributed by arXiv under the Creative Commons Attribution 4.0 International licence, http://creativecommons.org/licenses/by/4.0/, as recorded in the arXiv licence metadata for this exact version and shown on the abstract page https://arxiv.org/abs/2604.22705v1. Full licence notice: \texttt{licenses/arxiv-2604.22705-CC-BY-4.0.txt}.\par\smallskip

\noindent\textit{Editorial scope.} Counted unit: the article's lemma lem:wallpaper (source0.tex lines 336--338). The article's complete original proof of it is delivered with it (source0.tex lines 339--347, one \texttt{proof} environment of 9 lines). No mathematics is written, repaired or re-derived: the statement and the proof are the article's own lines, in the article's own notation, and no retained line is substituted or re-flowed.  No further source-local context is needed: every cross-reference of every retained line resolves inside the two delivered spans.  Nothing was omitted from any retained span, so the package carries no omission record.  No retained line contains a citation, so the wrapper carries no bibliography and no bibliography entry.  No retained line contains a figure environment, an image-inclusion command or drawing code, so no image is delivered and the package declares no asset dependency.  Editorial formatting only, in addition: an AMS \texttt{amsart} preamble that restates the article's own declarations and macros used by the retained lines, namely the environments \texttt{lemma} on separate counters, together with the standard packages those lines need.  The retained environments print in this document as Lemma 1 in order of appearance, whereas the article prints the same items with the numbering of its own full document.  The complete native source of the article as distributed in the arXiv e-print for this exact version is bundled unchanged as \texttt{sources/199118/source0.tex}.
\begin{lemma}\label{lem:wallpaper}
  Let $M$ be a locally-finite map with 1-end in $\EE^2$ such that $V(M)$ is a locally finite collection of vertices, and suppose $G\leq \text{Isom}(\EE^2)$ acts quasi-transitively on $M$. Then $G$ is a wallpaper group.
\end{lemma}

\begin{proof}
    It suffices to show that there exist linearly independent translations $t_1,t_2\in G$. We first suppose that $G$ contains no translations. Then all orbits $Gx$, $x\in\EE^2$, are finite. Since $G$ acts quasi-transitively this implies that $V(M)$ is finite, contradicting that $M$ has 1-end.\\
    Suppose that there is a single translation $t\in G$ up to linear dependence. Without loss of generality suppose that $t$ attains the minimum translation distance. Then $G$ is a Frieze group and has a finite index infinite-cyclic translation subgroup $T=\langle t\rangle$. Then $T$ has a fundamental domain $F$ that is the interior of two parallel lines $\ell_1,\ell_2\subset \EE^2$. Let $d=\text{dist}(\ell_1,\ell_2)$, this corresponds with the translation distance of $t$. 
    The set of distances between adjacent vertices $\mathcal{E}:=\{|u-v|:\{u,v\}\in E(\Gamma)\}$ must be finite so there exists a maximum $\max\{\mathcal{E}\}$.
    Choose $n>0$ sufficiently large that the translation distance of $nt$ is at least $\max\{\mathcal{E}\}$. Then $T_n=\langle nt\rangle$ is index $n$ in $T$ and has a fundamental domain $F_n$  so that $\Bar{F}_n$  consists of the connected union of $n$ copies of $\bar{F}$. It follows that $T_n$ is finite index in $G$. Since $G$ acts quasi-transitively the set $\bar{F}\cap V(M)$ is finite, and hence $\bar{F}_n\cap V(M)$ is finite. Then removing $\bar{F}_n$ from the plane leaves the disjoint union
     $$\EE^2-F_n=P_1\sqcup P_2,$$
    where $P_1$ and $P_1$ are infinite connected components of the plane. Moreover by construction we have that if $u\in V(M)\cap P_1$ and $v\in V(M)\cap P_2$ then they cannot be connected by an edge of $M$.
    Since the orbit of $F_n$ under $T_n$ is bi-infinite, it follows that the induced subgraphs on $P_1\cap V(M)$ and $P_2\cap V(M)$ are infinite and connected. Then $F_n\cap V(M)$ serves as a finite cut set leaving two infinite connected components, contradicting that $M$ has 1-end.
    \end{proof}

\end{document}
